\section{Trabajos futuros}
\label{sec:futuros}

Los trabajos pendientes se agrupan en cinco frentes, ordenados aproximadamente por
prioridad y por dependencia mutua.

\subsection{Consolidaci\'on del hardware: de placas de desarrollo a PCB propia}

El sistema opera hoy sobre placas de desarrollo (CY8CKIT-059 y ESP32~DevKitC) cableadas
entre s\'i. El 1 de agosto se dio el primer paso de la migraci\'on al adoptar una
\textbf{placa portadora intermedia} con asignaci\'on de pines definitiva y verificada
autom\'aticamente. Los pasos restantes son:

\begin{itemize}[itemsep=0.15em]
    \item Trasladar el dise\~no anal\'ogico completo a \textbf{KiCad}, aprovechando que los
    valores de componentes ya fueron recuperados del formato binario propietario y
    documentados con su nivel de confianza.
    \item Dise\~nar, fabricar y ensamblar la PCB del nodo ge\'ofono, integrando PSoC, ESP32,
    alimentaci\'on, conector de ge\'ofono y ranura de tarjeta SD.
    \item Resolver la alimentaci\'on aut\'onoma del nodo: dimensionamiento de bater\'ia,
    regulaci\'on de bajo ruido y estimaci\'on de autonom\'ia para una jornada de campo.
    \item Dise\~nar el encapsulado y el sistema de acople al terreno.
\end{itemize}

\subsection{Fuente s\'ismica repetible: el martillo de leva}

La calidad de la imagen de dispersi\'on depende directamente de la repetibilidad de la
fuente, y una maza operada manualmente introduce variabilidad en energ\'ia, punto de
impacto y verticalidad. Wood y Cox~\cite{Wood2012} cuantificaron el efecto del tipo de
fuente sobre el sesgo y la incertidumbre de la curva de dispersi\'on, lo que justifica el
esfuerzo de mecanizar el impacto.

El dise\~no del \textbf{martillo de leva} arranc\'o el 1 de agosto con dos modelos
paralelos: uno anal\'itico bidimensional en Python y uno tridimensional en Simscape
Multibody derivado del CAD. Ambos detectaron errores propios que est\'an pendientes de
resoluci\'on: una \textbf{interferencia axial entre la leva y el brazo} (que invalida los
resultados de la simulaci\'on 3D en su estado actual) y una \textbf{violaci\'on de la
conservaci\'on de energ\'ia} en el modelo 2D. Las tareas siguientes son corregir ambos
modelos, converger a un perfil de leva, construir el mecanismo y caracterizar
experimentalmente la repetibilidad de la energ\'ia entregada y del instante de impacto.

\subsection{Conectividad y flujo de datos en campo}

El flujo actual requiere que un navegador est\'e escuchando en el momento de la captura,
porque el maestro no persiste los datos: \'estos se espejan en vivo por WebSocket y el
navegador los acumula. Adem\'as, mientras el tel\'efono est\'a conectado al punto de acceso
del maestro, el sistema operativo m\'ovil tiende a cortar o despriorizar los datos
m\'oviles al no detectar salida a internet, lo que bloquea cualquier sincronizaci\'on
autom\'atica.

Est\'a dise\~nado y con firmware implementado ---aunque a\'un sin probar en hardware--- un
\textbf{modo ENLACE} del maestro que resuelve ambos problemas permiti\'endole asociarse a
una red con salida a internet mientras mantiene su punto de acceso. Quedan pendientes:

\begin{itemize}[itemsep=0.15em]
    \item Probar el modo ENLACE en hardware y verificar que los esclavos adopten
    correctamente el canal del maestro cuando \'este cambia al asociarse a un router.
    \item Corregir un bloqueo mutuo detectado en el escaneo de redes cuando existe un
    identificador de red guardado que no est\'a presente.
    \item Implementar el servidor de datos del lado del anfitri\'on para persistencia
    autom\'atica, eliminando la dependencia de un navegador activo durante la captura.
\end{itemize}

\subsection{Ampliaci\'on del arreglo y campa\~na de validaci\'on}

\begin{itemize}[itemsep=0.15em]
    \item Escalar el arreglo desde la configuraci\'on actual hasta el n\'umero de canales
    objetivo, y caracterizar experimentalmente el \textbf{apareamiento entre canales} y
    el \emph{jitter} de sincronizaci\'on real sobre el arreglo completo ---magnitud que
    hasta ahora s\'olo se acot\'o por presupuesto te\'orico y por medici\'on sobre pocos nodos.
    \item Implementar y ensayar las \textbf{configuraciones matriciales} de medici\'on
    previstas en la propuesta, m\'as all\'a del arreglo lineal.
    \item Ejecutar una campa\~na de medici\'on en un sitio con \textbf{informaci\'on
    geot\'ecnica de referencia disponible} (SPT o equivalente).
    \item \textbf{Contrastar cuantitativamente} los perfiles $V_S$ obtenidos por MASW
    contra esa referencia. Este es el criterio de \'exito \'ultimo del proyecto y la
    validaci\'on que a\'un no se realiz\'o.
\end{itemize}

\subsection{Caracterizaci\'on complementaria de la instrumentaci\'on}

\begin{itemize}[itemsep=0.15em]
    \item \textbf{Validaci\'on estad\'istica sobre m\'as unidades}: el estudio del DA-AFE se
    limit\'o a dos placas; extenderlo a una poblaci\'on mayor permitir\'ia afirmaciones a
    nivel poblacional.
    \item \textbf{Deriva de largo plazo y comportamiento t\'ermico}: caracterizar la
    estabilidad de la calibraci\'on a lo largo de una jornada de campo y frente a
    variaciones de temperatura, no evaluadas hasta ahora.
    \item \textbf{Optimizaci\'on de par\'ametros del controlador}: el diagn\'ostico con banda
    muerta m\'as estrecha mostr\'o margen de mejora; una resoluci\'on de ajuste m\'as fina
    podr\'ia lograrse con DAC de salida de corriente o topolog\'ias de inyecci\'on con
    amortiguador, sin modificar el m\'etodo de calibraci\'on.
    \item \textbf{Extensi\'on del DA-AFE a otras funciones}: la calibraci\'on de
    \emph{offset} de continua es una demostraci\'on de factibilidad; el mismo principio de
    asistencia digital no intrusiva puede aplicarse a otras tareas de calibraci\'on y
    optimizaci\'on anal\'ogica.
    \item \textbf{Caracterizaci\'on de ruido del sistema completo} referido a la entrada,
    y comparaci\'on con el l\'imite de ruido t\'ermico del ge\'ofono.
\end{itemize}

\subsection{Mejoras identificadas y documentadas, de menor prioridad}

Durante la auditor\'ia pre-campo se documentaron mejoras que no bloquean la operaci\'on:
prelectura de bloque en el PSoC para acelerar la descarga aproximadamente al doble;
posibilidad de reintentar la descarga desde la interfaz tras una detenci\'on (los datos
permanecen en la SD); y correcci\'on de dos detalles cosm\'eticos en la presentaci\'on de
confirmaciones de comando y de la vista previa de decimaci\'on.

\nota{Este es el bloque que la mesa evaluadora suele mirar con m\'as atenci\'on en una
Primera Presentaci\'on, porque define el alcance restante. Para la versi\'on de 15 p\'aginas
sugiero conservarlo casi completo (~1,5 p\'aginas), recortando \S7.6 y comprimiendo
\S7.5. Conviene adem\'as acordar con el tutor un \textbf{cronograma con fechas} para estos
frentes, que hoy no existe y que la mesa probablemente pida.}
